\section{Introduction}

Every year Israeli high-school students sit the \emph{Bagrut}, the national
matriculation examinations that gate university admission. A long-standing
question in educational policy is how much a student's outcome is fixed by the
socioeconomic status of the \emph{municipality} their school sits in, versus what
the \emph{school itself} does with its resources. This project answers that
question with a reproducible, cross-validated machine-learning pipeline.

\textbf{Main question.} Can municipal socioeconomic status and school-level
institutional resources, \emph{combined}, meaningfully predict Israeli
high-school Bagrut achievement and advanced-track participation? \textbf{Secondary
questions.} (A)~Which subject (Mathematics or English) is more \emph{resilient}
to socioeconomic and institutional disparity, and do schools exist that
outperform their peers despite limited resources? (B)~Which specific school-level
factors (budget allocation, class size, sector, or supervision type) most
strongly predict outcomes, \emph{independent} of municipal wealth?

We model four school-level targets: the average Bagrut grade and the advanced
(5-unit) participation rate, each for Mathematics and English. A municipal score
cannot capture what happens inside a specific school, and separating the two is
what makes the policy question answerable.

\textbf{Revisions since the interim stage.} The study began as a municipal-only
design. Following review we added the Ministry's school-level dataset and rebuilt
the pipeline from ingestion onward. Later revisions introduced cross-validation
grouped by school, per-target feature selection, explicit handling of
privacy-suppressed grades, a municipal-versus-school ablation, and a nested
protocol in which feature selection and tuning run inside training folds only.

\textbf{Code availability.} All data-processing and modelling code, with
per-stage documentation, is openly available at \href{\repourl}{\repourl}. Every
figure and number below traces to a file in that repository.
